{\Large\textbf{Superconducting mesh}}

Consider a mesh made from a flat superconducting sheet by drilling a dense grid
of small holes into it. Initially the sheet is in a non-superconducting state,
and a magnetic dipole of dipole moment $m$ is at a distance $a$ from the mesh
pointing perpendicularly towards the mesh. Now the mesh is cooled so that it
becomes superconducting. Next, the dipole is displaced perpendicularly to the
surface of the mesh so that its new distance from the mesh is $b$. Find the
force between the mesh and the dipole. The pitch of the grid of holes is much
smaller than both $a$ and $b$, and the linear size of the sheet is much larger
than both $a$ and $b$.
